% Part: first-order-logic
% Chapter: introduction
% Section: satisfaction

\documentclass[../../../include/open-logic-section]{subfiles}

\begin{document}

\olfileid{fol}{int}{sat}

\olsection{满足}

一旦知道了公式是什么，我们就可以直接进入一阶逻辑的语义学：
这里，基本的定义是\emph{结构}的定义。对于我们这里的简单语言，
一个结构~$\Struct{M}$ 只有三个组成部分：
一个称为\emph{论域}的非空集合 $\Domain{M}$，$\Obj a$ 在~$\Struct{M}$ 中所指派的对象，
以及 $\Obj P$ 在~$\Struct{M}$ 中对其为真的那些对象。
由~$\Obj a$ 所指派的对象记作~$\Assign{\Obj a}{M}$，而 $\Obj P$ 对其为真的那些对象组成的集合记作~$\Assign{\Obj P}{M}$。
一个结构~$\Struct{M}$ 就由这三者组成：$\Domain{M}$、$\Assign{\Obj a}{M} \in
\Domain{M}$ 和 $\Assign{\Obj P}{M} \subseteq \Domain{M}$。
一般情况会更复杂，因为会有许多谓词符号和常项符号，常项符号可以不止一位，并且还会有函数符号。

这已足够为不包含变元的公式给出满足的定义了。
其思路是给出一个与我们定义公式的方式相对应的归纳定义。
我们首先规定原子公式何时在~$\Struct{M}$ 中被满足，然后在此基础上规定，例如，
$\lnot !A$ 何时在~$\Struct{M}$ 中被满足（取决于 $!A$ 是否在~$\Struct{M}$ 中被满足）。
例如，我们可以定义：

\begin{enumerate}
  \item $\Atom{\Obj P}{\Obj a}$ 在~$\Struct{M}$ 中被满足，当且仅当
  $\Assign{\Obj a}{M} \in \Assign{\Obj P}{M}$。
  \item $\lnot !A$ 在~$\Struct{M}$ 中被满足，当且仅当 $!A$ 不在~$\Struct{M}$ 中被满足。
  \item $(!A \land !B)$ 在~$\Struct{M}$ 中被满足，当且仅当 $!A$ 在~$\Struct{M}$ 中被满足，并且 $!B$ 同样在~$\Struct{M}$ 中被满足。
\end{enumerate}

假设$\Domain{M} = \{0, 1, 2\}$，$\Assign{\Obj a}{M} = 1$，
且$\Assign{\Obj P}{M} = \{1, 2\}$。这个定义将告诉我们
$\Atom{\Obj P}{\Obj a}$ 在~$\Struct{M}$ 中被满足（因为 $\Assign{\Obj a}{M} =  1 \in \{1,2\} = \Assign{\Obj P}{M}$）。它进一步告诉我们 $\lnot \Atom{\Obj P}{\Obj a}$ 不在~$\Struct{M}$ 中被满足，进而 $\lnot\lnot \Atom{\Obj P}{\Obj
a}$ 被满足而 $(\lnot \Atom{\Obj P}{\Obj a} \land \Atom{\Obj P}{\Obj a})$ 不被满足，依此类推。

在为量词下定义时，麻烦出现了：
我们想说“$\lexists[\Obj
v_0][\Atom{\Obj P}{\Obj v_0}]$ 被满足当且仅当 $\Atom{\Obj P}{\Obj
v_0}$ 被满足。”但结构~$\Struct{M}$ 并没有告诉我们如何处理变元。
我们实际想说的是，\emph{对于~$\Obj v_0$ 的某个值而言}，$\Atom{\Obj P}{\Obj v_0}$ 被满足。
为了精确表达这一点，我们需要一种方法，将~$\Domain{M}$ 的元素不仅指派给 $\Obj a$，也指派给~$\Obj v_0$。
为此，我们引入变元\emph{指派}。
变元指派就是一个函数~$s$，它将变元映射到~$\Domain{M}$ 的元素（在我们的例子中，映射到 $1$、$2$ 或~$3$ 之一）。
由于我们无法预先知道哪些变元可能出现在公式中，所以不能限制 $s$ 仅给哪些变元赋值。
简单的解决办法是要求 $s$ 为\emph{所有}变元~$\Obj v_0$, $\Obj v_1$, \dots 赋值。我们只需用到其中所需的那部分。

我们将不再仅相对于一个结构来定义公式的满足，而是相对于
一个结构~$\Struct{M}$ \emph{以及}一个变元指派~$s$ 来定义它，
并简记为$\Sat{M}{!A}[s]$。现在我们的定义将包含一条额外的条款来处理含有变元的原子公式：

\begin{enumerate}
  \item $\Sat{M}{\Atom{\Obj P}{\Obj a}}[s]$ 当且仅当
  $\Assign{\Obj a}{M} \in \Assign{\Obj P}{M}$。
  \item $\Sat{M}{\Atom{\Obj P}{\Obj v_i}}[s]$ 当且仅当
  $s(\Obj v_i) \in \Assign{\Obj P}{M}$。
  \item $\Sat{M}{\lnot !A}[s]$ 当且仅当 非$\Sat{M}{!A}[s]$。
  \item $\Sat{M}{(!A \land !B)}[s]$ 当且仅当 $\Sat{M}{!A}[s]$ 且 $\Sat{M}{!B}[s]$。
\end{enumerate}

好，这解决了一个问题：我们现在可以说出 $\Struct{M}$ 何时对值~$s(\Obj v_0)$ 满足 $\Atom{\Obj P}{\Obj v_0}$。
为了正确定义 $\lexists[\Obj v_0][\Atom{\Obj P}{\Obj
v_0}]$，我们还需要做一件事：
我们希望：$\Sat{M}{\lexists[\Obj v_0][\Atom{\Obj P}{\Obj v_0}]}[s]$ 当且仅当
对于\emph{某种}给~$\Obj v_0$ 赋值的方式 $s'$，$\Sat{M}{\Atom{\Obj P}{\Obj v_0}}[s']$。
但指派给~$\Obj v_0$ 的值不一定非得是 $s(\Obj v_0)$ 所指派的那个值。
为此我们引入一个记号：如果 $m \in \Domain{M}$，那么我们令 $\Subst{s}{m}{\Obj v_0}$ 表示这样一个指派：它除了给 $\Obj v_0$ 指派~$m$ 之外，对其他所有变元的指派都与~$s$ 相同。现在我们的定义可以是：

\begin{enumerate}\setcounter{enumi}{4}
  \item $\Sat{M}{\lexists[\Obj v_i][!A]}[s]$ 当且仅当
  对某个~$m \in \Domain{M}$，$\Sat{M}{!A}[\Subst{s}{m}{\Obj v_i}]$。
\end{enumerate}

这样行得通吗？假设对于所有~$i \in \Nat$，我们令 $s(\Obj v_i) = 0$。
$\Sat{M}{\lexists[\Obj v_0][\Atom{\Obj P}{\Obj v_0}]}[s]$ 当且仅当
存在 $m \in \Domain{M}$ 使得 $\Sat{M}{\Atom{\Obj P}{\Obj
v_0}}[\Subst{s}{m}{\Obj v_0}]$。而这样的 $m$ 确实存在：我们可以选择 $m = 1$ 或 $m = 2$。
注意，即使由 $s$ 本身指派给~$\Obj v_0$ 的值——在此例中是 $0$——无法使其成立，上式依然成立。
我们有 $\Sat{M}{\Atom{\Obj P}{\Obj v_0}}[\Subst{s}{1}{\Obj
v_0}]$ 但没有 $\Sat{M}{\Atom{\Obj P}{\Obj v_0}}[s]$。

如果这看起来混乱且繁琐：确实如此。
但为了对所有公式的满足关系给出精确的归纳定义，这种额外的复杂性是必需的。
并且，为了精确定义语义概念，我们也需要类似的机制。
当然，还有其他方式可以做到这一点，但它们同样（不）优雅。

\end{document}